\section{Axiomas de separabilidad y numerabilidad}

% Capítulo 18: Primer axioma de numerabilidad

\subsection{Espacios primero numerables}

\begin{definición}[Primer axioma de numerabilidad]
    Un espacio topológico $X$ es \textbf{primero numerable} si todo punto $x\in X$ tiene una base local numerable.
\end{definición}

\ejemplo{Todo espacio métrico es primero numerable: en $x$ basta la familia de bolas $B(x,1/n)$, $n\in\mathbb{N}$.}

\subsection{Sucesiones y convergencia}

\begin{proposición}
    En un espacio primero numerable, un punto $x$ pertenece a $\mathrm{Cl}(A)$ si y sólo si existe una sucesión $(a_n)$ en $A$ tal que $a_n\to x$.
\end{proposición}

\subsection{Continuidad secuencial}

\begin{proposición}
    Si $X$ es primero numerable, para $f:X\to Y$ se tiene:
    $$f \text{ continua } \Longleftrightarrow \forall (x_n),\ x_n\to x \Rightarrow f(x_n)\to f(x).$$
\end{proposición}

% Capítulo 19 (Axiomas): Segundo axioma de numerabilidad

\subsection{Espacios segundo numerables}

\begin{definición}[Segundo numerable]
    Un espacio $X$ es \textbf{segundo numerable} (o verifica el \emph{segundo axioma de numerabilidad}) si admite una base numerable para su topología.
\end{definición}

\begin{proposición}
    Todo espacio 2-numerable es 1-numerable.
\end{proposición}
\begin{proof}
    Si $\mathcal{B}$ es una base numerable, para cada $x \in X$ la colección $\mathcal{B}_x = \{B \in \mathcal{B} : x \in B\}$ es una base local numerable (por ser subconjunto de $\mathcal{B}$).
\end{proof}

\ejemplo{
    \begin{itemize}
        \item La topología usual de $\mathbb{R}^n$ es 2-numerable: los intervalos $(q_1 - r, q_1 + r) \times \cdots$ con centros racionales y radio racional forman una base numerable.
        \item La recta de Sorgenfrey es 1-numerable pero \textbf{no} 2-numerable: si tuviese base numerable $\mathcal{B}$, la asignación $x \mapsto B_x$ (donde $x \in B_x \subset [x, x+1)$) sería inyectiva $\mathbb{R} \to \mathcal{B}$, pero $\mathbb{R}$ no es numerable.
    \end{itemize}
}

\subsection{Espacios separables}

\begin{definición}[Denso y separable]
    Un subespacio $A \subset X$ es \textbf{denso} en $X$ si $\mathrm{Cl}(A) = X$, equivalentemente si $A$ corta a todos los abiertos no vacíos.

    Un espacio topológico $X$ es \textbf{separable} si contiene un subespacio denso y numerable.
\end{definición}

\ejemplo{
    $\mathbb{Q}$ es denso y numerable en $\mathbb{R}$ con la topología usual (y también en la Sorgenfrey), luego $\mathbb{R}$ es separable. Asimismo, $\mathbb{R}^n$ es separable tomando $\mathbb{Q}^n$ como subconjunto denso numerable.
}

\begin{proposición}
    Si un espacio topológico es 2-numerable, entonces es separable.
\end{proposición}
\begin{proof}
    Sea $\mathcal{B}$ una base numerable de $\tau$. En cada abierto básico $B \in \mathcal{B}$ escogemos un punto $x_B$. El conjunto $A = \{x_B : B \in \mathcal{B}\}$ es numerable y denso: si $U$ es un abierto no vacío, existe $B$ básico con $x \in B \subset U$, luego $x_B \in A \cap U \neq \emptyset$.
\end{proof}

\begin{proposición}
    Sea $(X, \tau)$ un espacio topológico metrizable. Si $X$ es separable, entonces es 2-numerable.
\end{proposición}
\begin{proof}
    Sea $A$ un subespacio denso y numerable con métrica $d$. Las bolas $B(a, r)$ con $a \in A$ y $r \in \mathbb{Q}_{>0}$ forman una base numerable: dado $x \in U$ abierto, existe $B(x, \varepsilon) \subset U$; tomando $a \in A \cap B(x, \varepsilon/4)$ y racional $\varepsilon/4 < r < 3\varepsilon/4$ se tiene $x \in B(a, r) \subset B(x, \varepsilon) \subset U$.
\end{proof}

\begin{corolario}
    La topología de Sorgenfrey en $\mathbb{R}$ no procede de ninguna métrica: es separable (pues $\mathbb{Q}$ es denso) pero no es 2-numerable.
\end{corolario}

% Capítulo 20: Propiedades de separación

\subsection{Axiomas de separación}

\begin{definición}[$T_0$, $T_1$, Hausdorff]
    Un espacio $X$ es:
    \begin{itemize}
        \item $T_0$ si para $x\neq y$ existe un abierto que contiene a uno y no al otro.
        \item $T_1$ si para $x\neq y$ existen abiertos que separan cada punto del otro.
        \item \textbf{Hausdorff} ($T_2$) si para $x\neq y$ existen abiertos disjuntos $U,V$ con $x\in U$, $y\in V$.
    \end{itemize}
\end{definición}

\begin{proposición}
    En un espacio Hausdorff, los límites de sucesiones (cuando existen) son únicos.
\end{proposición}

\begin{proposición}
    Todo subespacio de un Hausdorff es Hausdorff y todo producto finito de Hausdorff es Hausdorff.
\end{proposición}

\subsection{Espacios normales}

\begin{definición}[Regular y normal]
    Un espacio $X$ es regular si puntos y cerrados disjuntos pueden separarse por abiertos. Es normal si cualesquiera dos cerrados disjuntos pueden separarse por abiertos.
\end{definición}

\begin{proposición}
    Todo espacio métrico es normal (y en particular regular y Hausdorff).
\end{proposición}
